\documentclass[11pt]{article}
\usepackage[margin=1in]{geometry}
\usepackage{booktabs}
\usepackage{hyperref}
\usepackage{longtable}
\usepackage{parskip}

\title{Replication Report: Br\"uggemann et al.\ (2018)\\
\large Pan-genome analysis of the genus \emph{Finegoldia} identifies two distinct clades,\\
strain-specific heterogeneity, and putative virulence factors}
\author{Ollie (OpenClaw AI) --- Replication Wave\\BVBRC-39 (TOPUP85 rank-22)}
\date{2026-07-01}

\begin{document}
\maketitle

\section*{Front matter}
\textbf{Paper:} Br\"uggemann H, Jensen A, Nazipi S, Aslan H, Meyer RL, Poehlein A, Brzuszkiewicz E, Al-Zeer MA, Brinkmann V, S\"oderquist B. \emph{Scientific Reports} 8:266 (2018).\\
\textbf{DOI:} \href{https://doi.org/10.1038/s41598-017-18661-8}{10.1038/s41598-017-18661-8} \quad
\textbf{PMC:} PMC5762925 \quad
\textbf{PMID:} 29321635 \quad
\textbf{Open access:} CC BY 4.0.\\
\textbf{Verdict:} \textbf{REPLICATED} (same 17 genomes, independent tools; all central quantitative claims reproduced to within $\sim$1\%).

\section{Paper summary}
The authors sequenced 10 new \emph{Finegoldia} isolates (Sweden) and combined them with 7 previously published \emph{F.\ magna} genomes (\textbf{17 total}). Using core-genome SNP phylogeny (Parsnp), ANI (JSpeciesWS), and pan-genome analysis (ProteinOrtho), they found that the genus splits into \textbf{two distinct clades} separated by an ANI of \textbf{90.7\%}: one clade is \emph{F.\ magna sensu stricto}, the other a more heterogeneous, more abundant novel species tentatively named \textbf{``\emph{Finegoldia nericia}.''} They report a \textbf{core proteome of 1202 orthologs} shared by all 12 analyzed genomes ($\sim$68\% of each strain's CDS), extensive strain-specific heterogeneity in the accessory genome, a conserved sortase-dependent pilus locus (Fmp1/Fmp2 + sortases), \textbf{2--4 CAMP-factor copies} per genome, and heterogeneous distribution of host-interacting virulence factors (protein L, PAB, FAF, SufA) --- notably protein L present in only $\sim$10\% of isolates.

\section{Claims tested}
\begin{longtable}{clllc}
\toprule
\# & Claim & Type & Testable? & Tested? \\
\midrule
C1 & The 17 genomes are publicly available. & Data availability & Yes & \checkmark \\
C2 & $\sim$1570--1906 CDS/genome (avg 1760); GC $\sim$32\%; sizes $\sim$1.7--2.0 Mb. & Genomic stats & Yes & \checkmark \\
C3 & Population splits into 2 clades, inter-clade ANI 90.7\%. & Phylogenomic/ANI & Yes & \checkmark \\
C4 & 12-genome subset = 4 \emph{magna} + 8 ``\emph{nericia}.'' & Clustering & Yes & \checkmark \\
C5 & Core proteome = 1202 orthologs ($\sim$68\% of avg CDS). & Pan-genome & Yes & \checkmark \\
C6 & Strain-specific heterogeneity / open accessory genome. & Pan-genome & Yes & \checkmark \\
C7 & Sortase-dependent pilus locus (Fmp1/Fmp2 + sortases) conserved. & Genomic & Yes (partial) & \checkmark \\
C8 & 2--4 CAMP-factor copies per genome. & Genomic & Yes & \checkmark \\
C9 & Heterogeneous virulence-factor distribution; protein L in $\sim$10\%. & Genomic & Yes & \checkmark \\
\bottomrule
\end{longtable}

\section{Method}
Same 17 genomes as the paper; independent tools. The paper's WGS accessions were mapped 1:1 to current NCBI GCA assemblies (by both WGS-project prefix and strain name), so this is a \emph{same-data} replication, not a proxy.

\begin{enumerate}
\item \textbf{Full text} pulled from Europe PMC REST (\texttt{/PMC5762925/fullTextXML}) and parsed for accessions and claims.
\item \textbf{Accession mapping} --- queried NCBI Datasets taxon report (278 \emph{Finegoldia} records); matched all 17 paper strains to GCA accessions (\texttt{work/paper\_17\_map.tsv}).
\item \textbf{Download} --- \texttt{datasets download genome accession --inputfile acc\_list.txt --include genome,protein,gff3}.
\item \textbf{Genome stats} (\texttt{genome\_stats.py}) --- length, contigs, GC\%, CDS count from PGAP \texttt{protein.faa}.
\item \textbf{ANI / clades} (\texttt{ani\_analysis.py}, \texttt{ani\_cluster2.py}) --- \textbf{fastANI} all-vs-all (17$\times$17); distance = 100$-$ANI; SciPy average-linkage cut into 2 clusters; clades named by presence of \emph{F.\ magna} type strain ATCC 29328.
\item \textbf{Pan/core genome} (\texttt{pangenome.py}) --- concatenated 12-genome subset proteomes, clustered with \textbf{CD-HIT} (c=0.5, n=3); families present in all 12 = core; families in 1 strain = singletons.
\item \textbf{Virulence factors} (\texttt{vf\_survey.py}) --- curated UniProt \emph{F.\ magna} reference proteins (FAF, SufA, sortase, CAMP factor, PAB, protein L, albumin-binding homolog); \textbf{blastp} vs each proteome; presence rule pident $\geq$ 40 and coverage $\geq$ 50\%; CAMP paralog count at relaxed pident $\geq$ 30 and cov $\geq$ 40\%.
\item \textbf{LLM-judge} --- Argo free \texttt{argo:gpt-5.2} (opus-4.8 fallback) scored coverage/agreement/verdict from claims + results.
\end{enumerate}

\section{Results vs paper}

\subsection{Genome statistics (C2) --- near-exact}
\begin{tabular}{lll}
\toprule
Metric & Paper & This replication (NCBI PGAP) \\
\midrule
CDS/genome & 1570--1906, avg 1760 & 1563--1956, mean 1759 \\
GC\%       & $\sim$32\%           & 31.7--32.1\% \\
Genome size & $\sim$1.7--2.0 Mb   & 1.68--2.03 Mb \\
\bottomrule
\end{tabular}

\subsection{Two clades and ANI (C3) --- reproduced}
\begin{tabular}{lll}
\toprule
Quantity & Paper & This replication (fastANI) \\
\midrule
Number of clades & 2 & \textbf{2} \\
Inter-clade ANI  & \textbf{90.7\%} & 90.67--91.70\% (min 90.67) \\
Intra-clade ANI  & (high, within-species) & mean 96.06\% \\
Clade sizes (17 total) & magna vs nericia & magna = 9, nericia = 8 \\
\bottomrule
\end{tabular}

The minimum inter-clade ANI (90.67\%) essentially lands on the paper's reported 90.7\% --- an independent ANI engine (fastANI vs the paper's JSpeciesWS) recovered the same species boundary.

\subsection{12-genome subset split (C4) --- exact}
Paper: 4 \emph{magna} + 8 \emph{nericia}. This replication: 4 \emph{magna} + 8 \emph{nericia} (exact match).
\begin{itemize}
\item \emph{magna} (4): 07T609, 08T492, 09T408, ATCC 29328
\item \emph{nericia} (8): 09T494, 12T272, 12T273, 12T306, ATCC 53516, CCUG 54800, T151023, T160124
\end{itemize}
Notably, two strains historically deposited as ``\emph{F.\ magna}'' (\textbf{ATCC 53516} and \textbf{CCUG 54800}) fall in the \emph{nericia} clade --- precisely the paper's central point.

\subsection{Core / pan-genome (C5, C6) --- near-exact}
\begin{tabular}{llll}
\toprule
Quantity & Paper (ProteinOrtho) & This replication (CD-HIT) & Match \\
\midrule
Core orthologs (all 12) & \textbf{1202} & \textbf{1209} & +0.6\% \\
Core as \% of avg CDS   & 68\%          & 69.9\%        & \checkmark \\
Pan-genome families     & (open)        & 2992          & --- \\
Singletons (1 strain)   & ``many''      & \textbf{892}  & \checkmark \\
\bottomrule
\end{tabular}

Gene-frequency distribution ($k$ strains $\to$ \# families): 1$\to$892, 2$\to$222, 3$\to$169, \ldots, 11$\to$97, 12$\to$1209. The large singleton/cloud fraction directly confirms strain-specific heterogeneity (C6).

\subsection{Virulence factors (C7, C8, C9) --- confirmed}
CAMP-factor copy number (C8): all 17 genomes carry \textbf{2 copies} (within the paper's 2--4 range).

\begin{tabular}{lrrrrl}
\toprule
Virulence factor & /17 & \% & magna & nericia & Paper says \\
\midrule
CAMP factor              & 17 & 100\% & 9/9 & 8/8 & in all (2--4 copies) \\
SufA (subtilisin)        & 17 & 100\% & 9/9 & 8/8 & conserved host factor \\
FAF (adhesion)           & 12 & 70\%  & 8/9 & 4/8 & heterogeneous \\
Albumin-binding homolog  &  9 & 52\%  & 5/9 & 4/8 & heterogeneous \\
PAB (albumin-binding)    &  8 & 47\%  & 5/9 & 3/8 & heterogeneous \\
\textbf{Protein L (Ig-L)} & \textbf{2} & \textbf{11\%} & 0/9 & 2/8 & \textbf{$\sim$10\%} \\
\bottomrule
\end{tabular}

The protein-L result (11\%) matches the paper's headline heterogeneity figure ($\sim$10\%) almost exactly, using an independent homology search. Sortase genes were found in every genome (4--9 per genome by annotation), supporting the conserved-pilus-machinery claim (C7).

\subsection{Corpus availability}
NCBI Datasets now indexes \textbf{278 \emph{Finegoldia} records} (168 GCA primary) vs 17 in 2018 --- fully public and greatly expanded, enabling much larger future pan-genome studies.

\section{Verdict}
\textbf{REPLICATED.} All nine stated claims were tested on the same 17 genomes the paper analyzed, using independent tools (fastANI, CD-HIT, blastp) rather than the authors' stack (Parsnp, JSpeciesWS, ProteinOrtho, Prokka). Every central number reproduced to within $\sim$1\%: two clades at inter-clade ANI 90.67\% (paper 90.7\%); 4 \emph{magna} + 8 \emph{nericia} among the 12-genome subset (exact); core proteome 1209 (paper 1202); mean 1759 CDS/genome (paper 1760); protein L in 11\% of strains (paper $\sim$10\%). Independent LLM-judge (Argo free gpt-5.2): REPLICATED, coverage 9/9, agreement 9/9.

\section{Genuine critique}
This report reproduces the paper's headline numbers, but the strength of that agreement should be interpreted with the following honest caveats:

\begin{enumerate}
\item \textbf{Same input, different tools --- not a fully independent test.} The 17 assemblies are the same physical genome data the authors deposited; only the analytical stack differs. A ``deeper'' replication would resequence or reassemble at least a subset of the isolates from raw reads, which was not done here. Any systematic bias baked into NCBI PGAP annotation (CDS calls, hypothetical labeling) propagates into both the CDS-count and pan-genome numbers.
\item \textbf{Substituted analytical equivalents, not identical methods.} The paper's core phylogeny (Parsnp, 126{,}647 SNPs) and ProteinOrtho pan-genome were \emph{not} regenerated. fastANI + CD-HIT recover the same clade structure and a very similar core size, but they answer slightly different questions: fastANI is a coarse whole-genome ANI approximation, not a SNP-level alignment, and CD-HIT at c=0.5, n=3 is a permissive clustering that is known to slightly over-merge distantly related paralogs (which likely explains why our core count came in \emph{higher} at 1209 vs 1202 --- 0.6\% is within noise but the sign of the discrepancy is expected from the tool substitution, not luck).
\item \textbf{Virulence-factor calls are threshold-dependent.} The blastp presence rule (pident $\geq$ 40, coverage $\geq$ 50\%) was chosen a priori but is a soft cutoff. The protein-L call (2/17 = 11\%) is within one hit of matching the paper exactly; a slightly stricter threshold would have given 1/17 (5.9\%) and a slightly looser one 3/17 (17.6\%). The agreement is real but the exact number is not robust to $\pm$5\% pident. FAF (12/17) and albumin-binding (9/17) are similarly threshold-sensitive because the paper never quantified them numerically, only called them ``heterogeneous.''
\item \textbf{C7 (pilus locus) was only partially tested.} We confirmed sortase presence (4--9 per genome by annotation) and inferred that the machinery is conserved, but we did \emph{not} reconstruct the Fmp1/Fmp2 gene-cluster synteny, and PGAP labels many pilus-associated proteins ``hypothetical protein,'' which caps annotation-based inference. A proper test would run a synteny + phylogeny of the Fmp locus per strain.
\item \textbf{Clade assignment used a 2-way average-linkage cut of the ANI distance matrix} --- the number of clades was fixed \emph{a priori} to 2 to match the paper's claim. Single-linkage at the standard 95\% species threshold reveals finer sub-structure (3 groups). We are testing whether the paper's boundary is stable, not discovering the boundary de novo.
\item \textbf{LLM-judge is a single-judge, single-model scoring.} Argo free gpt-5.2 (with opus-4.8 fallback) returned 9/9 coverage and 9/9 agreement, but a single LLM is not an independent adjudicator --- it saw the same claims/results text as this report. It confirms internal self-consistency, not external validity. A multi-judge or human-adjudicated scoring would be stronger.
\item \textbf{Nothing was fabricated, but the CAMP-factor result is worth flagging as ``consistent but not strong.''} The paper says 2--4 copies; every one of our 17 genomes returned exactly 2. That is within the paper's stated range, but the modal value (2) may be a floor effect of our paralog-detection threshold (pident $\geq$ 30, cov $\geq$ 40\%) rather than an accurate count --- the paper likely used manual curation to distinguish 3rd/4th paralogs.
\item \textbf{We did not attempt the paper's BRIG comparative-genomics figures, mobile-element inventory, or CRISPR analysis.} These were secondary in the paper and are not part of the core claims tested (C1--C9), but they are unverified here.
\end{enumerate}

Bottom line: the verdict REPLICATED is well-supported for the paper's \emph{quantitative headline claims} (clades, ANI, core size, CDS/genome, protein-L rarity), but the reader should treat this as a \emph{same-data, tool-substitution replication with soft thresholds}, not a fully independent from-raw-reads reproduction.

\section{Open questions}
Five substantive open problems (with concrete next steps) are enumerated in the companion machine-readable file \texttt{open\_questions.json}, covering: (i) raw-read reassembly and the effect on CDS/GC/annotation, (ii) formal taxonomic validation of \emph{F.\ nericia}, (iii) full Fmp pilus locus synteny + phylogeny, (iv) revisiting clade structure using the 168 GCA assemblies now available (vs the original 17), and (v) mechanistic investigation of protein L's clade-specific distribution.

\section{Reproducibility artifacts}
See \texttt{workflow.md}, \texttt{artifacts\_summary.md}, and \texttt{failure\_analysis.md}. All inputs are free and public (Europe PMC, NCBI Datasets, UniProt); wall-clock $\sim$5 min on a laptop.

\end{document}
